\documentclass[a4paper,10pt]{article}
\usepackage[utf8]{inputenc}
\usepackage[T1]{fontenc}
\usepackage{lmodern,booktabs,longtable,geometry,hyperref,adjustbox}
\geometry{margin=2.3cm}
\setlength{\parindent}{0pt}\setlength{\parskip}{4pt}\emergencystretch=2em
\begin{document}
\begin{center}\Large\bfseries Seguimiento para una posible v1.5: resultados completos\end{center}

La versión de envío v1.4.2 permanece congelada. Este informe evalúa los
cuatro trabajos propuestos; no constituye una nueva versión enviada ni
una promesa de aceptación editorial.

\section*{Resultado principal y decisión}

No se establece una ventaja temporal general frente a L2 al exigir el objetivo común de precisión y calidad.
Budget-L2 alcanza 54/72 objetivos primarios, frente a 51/72 de L2.
La razón del indicador de disponibilidad verificada PAR2 es 0.876 [0.441, 1.722].
Entre los 44 éxitos comunes, la razón de disponibilidad real es 0.980.
PAR2 asigna 60 segundos de resolución más preparación a un objetivo no
alcanzado; su razón no es una reducción del tiempo observado de las soluciones.
En las cuatro especificaciones del experimento de precisión, L2 satisface todos los presupuestos finales comprobados. No se demuestra recuperación primal adicional frente a L2 en esa muestra.

El avance tangible del software es un selector que incorpora la exportación
efectiva a la decisión y un auditor independiente de sus reservas. Es una
extensión condicional del procedimiento existente, no una prueba de un
nuevo escalador universalmente competitivo. La información física conserva
su valor de especificación; las etiquetas no afectan al solver una vez
compilado el mismo vector de factores.

\section*{Diseño y significado de los tiempos}

630 asignaciones registradas: 288 para candidatos, 126 para precisión y
216 para robustez. Una repetición técnica predeterminada corrige la
implementación de tiempos de verificación del primer caso Base; no añade
un escenario ni una semilla. Se conservan un fallo previo al optimizador,
el programa original, la enmienda y todos los resultados iniciales.
La repetición solo sustituye tiempos comparativos si vectores, trabajo,
endpoints y primer objetivo se reproducen exactamente.

24 escenarios conocidos de 24, 72 y 168 horas; semillas nuevas 211, 257,
293; cuatro hilos, límite común de 30 segundos, gap nativo solicitado del 1\%.
CPLEX utiliza doce escenarios seleccionados por índice y búsqueda
tradicional común para instrumentar sus incumbentes. El control de
ScaleFlag usa seis escenarios: automático, 0 y 2. Son escenarios de la
misma familia y padres conocidos, no una muestra externa nueva.

El objetivo primario permite un costo hasta U\_ref + 0,001 max(1,|U\_ref|);
el secundario usa 0,01. U\_ref es un testigo factible racional adquirido
antes de estas ejecuciones, no un óptimo certificado. El gap de parada
del 1\% puede terminar una búsqueda sin llegar al objetivo primario del
0,1\%, incluso sin error numérico. Un objetivo no alcanzado no equivale
automáticamente a una cota incorrecta.

Se verifican presupuestos físicos crudos y tras redondear enteros,
distancia entera <=1e-9 y cotas exactas del punto redondeado. Las filas
sin unidad declarada usan error hacia atrás <=1e-6. Esta definición
no certifica factibilidad exacta de todas las filas ni optimalidad.
Las eventuales violaciones diminutas de cotas se conservan como tales.
El callback guarda todos los candidatos y el punto final; el verificador
corre después del solver. El tiempo de descubrimiento es retrospectivo;
la disponibilidad real incluye la espera hasta fin del solver y el
costo de comprobar el prefijo de candidatos. No se presenta como un
verificador en línea ni como un nuevo criterio de parada.

\section*{I. Objetivos comprobados y calidad}

\begin{center}\small\adjustbox{max width=\linewidth}{\begin{tabular}{llllll}\toprule
Política & Objetivo 0,1\% & Objetivo 1\% & Presupuestos finales & Dominios finales & Contradicciones B\\\midrule
Base & 53/72 & 68/72 & 72/72 & 69/72 & 2\\
Budget-L2 & 54/72 & 72/72 & 72/72 & 72/72 & 1\\
L2 & 51/72 & 66/72 & 72/72 & 72/72 & 7\\
L2-NF3 & 54/72 & 72/72 & 72/72 & 72/72 & 0\\
\bottomrule\end{tabular}}\end{center}

Las contradicciones usan testigos exactos de la representación efectivamente
importada; ausencia de contradicción no significa certificación de todas
las cotas. Un candidato que pasa el control primal puede coexistir con
una cota refutada.

\begin{center}\small\adjustbox{max width=\linewidth}{\begin{tabular}{lllll}\toprule
Control & Pares con éxito común & PAR2 descubrimiento [95\%] & PAR2 disponibilidad [95\%] & Disponibilidad en éxitos comunes\\\midrule
Base & 45 & 0.772 [0.375, 1.596] & 0.802 [0.405, 1.595] & 0.742\\
L2 & 44 & 0.875 [0.422, 1.807] & 0.876 [0.441, 1.722] & 0.980\\
L2-NF3 & 46 & 0.953 [0.497, 1.834] & 0.943 [0.509, 1.746] & 0.994\\
\bottomrule\end{tabular}}\end{center}

Los intervalos remuestrean 24 escenarios, conservando juntas las tres semillas:
100000 remuestras, percentiles marginales 95\%. Son condicionales a los
padres conocidos. Los éxitos comunes son un subconjunto seleccionado,
por eso se publican junto al análisis penalizado de todos los casos.
Al objetivo secundario, Budget-L2 alcanza 72/72 y L2 66/72;
su PAR2 de disponibilidad es 0.753 [0.454, 1.085].
Entre éxitos comunes la razón es 1.041.
El indicador penalizado puede mejorar por mayor cumplimiento del objetivo,
aunque los éxitos comunes no lleguen antes. No se confunden ambos efectos.

\section*{II. Precisión y necesidad del contrato}

Requisitos numéricos experimentales, sin atribuirles autoridad operativa:

\noindent\textbullet\ b10: potencia 0,01 MW; balance y activación 10 m3.\par
\noindent\textbullet\ b100: potencia 0,01 MW; balance y activación 100 m3.\par
\noindent\textbullet\ b1000: potencia 0,01 MW; balance y activación 1000 m3.\par
\noindent\textbullet\ heterogeneous: potencia 0,001 MW; balance 0,001 m3; activación 10 m3.\par

El balance a un litro es una meta de conservación numérica del estudio,
no una tolerancia de medición del operador. Las conversiones a unidades
algebraicas se reconstruyen desde cada fila. Cada celda tiene seis modelos
y tres semillas; los controles reutilizados no cuentan como nuevas llamadas.

\begin{center}\small\adjustbox{max width=\linewidth}{\begin{tabular}{lllll}\toprule
Contrato & Política & Presupuestos finales & Máximo q=v/delta & Filas que incumplen\\\midrule
b10 & Budget-L2 & 18/18 & 6.44e-08 & 0\\
b10 & Export-Budget-L2 & 18/18 & 7e-08 & 0\\
b10 & L2 & 18/18 & 6.44e-08 & 0\\
b100 & Budget-L2 & 18/18 & 6.77e-09 & 0\\
b100 & Export-Budget-L2 & 18/18 & 7e-09 & 0\\
b100 & L2 & 18/18 & 6.44e-09 & 0\\
b1000 & Budget-L2 & 18/18 & 6.44e-10 & 0\\
b1000 & Export-Budget-L2 & 18/18 & 6.34e-10 & 0\\
b1000 & L2 & 18/18 & 6.44e-10 & 0\\
heterogeneous & Budget-L2 & 18/18 & 0.000681 & 0\\
heterogeneous & Export-Budget-L2 & 18/18 & 0.000712 & 0\\
heterogeneous & L2 & 18/18 & 0.000644 & 0\\
\bottomrule\end{tabular}}\end{center}

Con activación de 1 m3, 22/24 modelos contienen 8256 filas
cuyo E=gamma sum |a\_entero| agota el presupuesto. Esto certifica incompatibilidad
de la condición suficiente, no infactibilidad del despacho ni imposibilidad
de que un punto concreto cumpla un metro cúbico. No se lanzan MIP en ese caso.
El control numérico compilado reproduce los factores antiguos sin diferencias (0 discrepancias).
Por tanto, no se atribuye un efecto causal a las etiquetas físicas. Los
márgenes por familia, percentiles de q, intervenciones y reservas R/delta
están íntegros en los CSV, incluyendo la aplicación de nuevos requisitos
a las soluciones antiguas Budget100.

\section*{III. CPLEX e interacción con el escalado interno}

\begin{center}\small\adjustbox{max width=\linewidth}{\begin{tabular}{llllll}\toprule
Configuración & Escenarios & Objetivo Budget & Objetivo L2 & PAR2 disponibilidad B/L2 [95\%] & B Budget/L2\\\midrule
CPLEX & 12 & 15/36 & 15/36 & 1.021 [1.005, 1.040] & 5/6\\
Gurobi ScaleFlag=-1 & 6 & 14/18 & 14/18 & 0.934 [0.128, 4.635] & 0/3\\
Gurobi ScaleFlag=0 & 6 & 15/18 & 13/18 & 0.689 [0.098, 3.571] & 0/4\\
Gurobi ScaleFlag=2 & 6 & 15/18 & 14/18 & 0.782 [0.109, 4.010] & 0/4\\
\bottomrule\end{tabular}}\end{center}

El criterio primario de dominios exige cotas exactas. En los puntos finales
de CPLEX, algunas desviaciones de cota están entre 1,3e-13 y 4,2e-11:
no se deben interpretar como una falta importante de precisión física.
Para comprobar su influencia se añadió una sensibilidad posterior: redondear
enteros, proyectar sobre las cotas exactas y volver a verificar el punto
completo, aplicando la misma regla a todas las políticas. No altera la
definición primaria ni mide una disponibilidad en línea del punto reparado.

\begin{center}\small\adjustbox{max width=\linewidth}{\begin{tabular}{llll}\toprule
Ámbito con proyección posterior & Política & Objetivo 0,1\% & Objetivo 1\%\\\midrule
I & Base & 53/72 & 71/72\\
I & Budget-L2 & 54/72 & 72/72\\
I & L2 & 51/72 & 66/72\\
I & L2-NF3 & 54/72 & 72/72\\
CPLEX & Budget-L2 & 30/36 & 36/36\\
CPLEX & Base & 24/36 & 36/36\\
CPLEX & L2 & 30/36 & 36/36\\
\bottomrule\end{tabular}}\end{center}

La proyección permite separar el efecto de cotas diminutamente excedidas
de la precisión de las filas y de la calidad del costo. No elimina las
contradicciones independientes de cotas duales ya detectadas. Las tablas
de tiempo primarias conservan su criterio registrado y no se reinterpretan
como tiempos medidos del posprocesamiento añadido.

Estas comparaciones usan las mismas representaciones y ajustes dentro de
cada solver. CPLEX se ejecutó con la biblioteca completa 22.1.2; la API
Community instalada por defecto no admite los tamaños mayores. No se
comparan sus ticks con las unidades de trabajo de Gurobi. Cambiar ScaleFlag
aísla una opción del solver, no reconstruye todas sus operaciones internas.
Los intervalos de seis o doce escenarios tienen incertidumbre limitada por
ese tamaño y no prueban generalidad entre familias.

\section*{Selector con exportación efectiva}

La auditoría independiente certifica 25344 casos fila-presupuesto del nuevo selector
(seis modelos y cuatro contratos; no son observaciones independientes)
con epsilon+B(alpha)+alpha E <= alpha delta, usando los coeficientes/RHS
realmente escritos. Si falta una cota finita y el producto tiene error
no nulo, ese candidato se rechaza. Un producto exacto no necesita esa cota.
Se elige entre candidatos finitos; abstenerse no prueba inexistencia global.
El orden logarítmico se calcula en binario64; las desigualdades residuales
se certifican con racionales exactos.

En la configuración de referencia, los factores del Budget-L2 anterior
también satisfacen la condición completa de exportación en todas las filas
físicas auditadas. Aquí el nuevo selector no recupera una garantía que
esos factores hubieran perdido; incorpora su comprobación a la elección.
Esta diferencia limita cuánto puede afirmarse de su utilidad adicional.

Se comprueban 2000 puntos racionales en la frontera de la tolerancia exportada.
Hay un ejemplo donde la proyección ideal no cumple la reserva uniforme
de exportación y un factor vecino sí la cumple. Es una prueba controlada
del algoritmo, no evidencia de una ventaja en una segunda aplicación.
La variante ordenada reproduce los 24 archivos completos del experimento II por SHA-256.
Medianas de selección: exhaustiva 8.175 s; ordenada 0.630 s.
Razón geométrica de tiempos de selección: 0.082.
Es un benchmark de compilación posterior; no reemplaza los tiempos de
proceso de los MIP ya registrados ni demuestra una aceleración completa.
La garantía sigue condicionada al residuo escrito, las cotas y la
distancia entera efectivamente satisfechos. No certifica optimalidad.

\section*{Qué incorporar al artículo}

La campaña responde las cuatro líneas solicitadas. La decisión editorial
debe basarse en cumplimiento y generalidad, no en escoger el cociente más
favorable. Este informe conserva todos los resultados y no modifica v1.4.2.
El selector ofrece una mejora concreta de implementación y certificación;
los resultados de candidatos y robustez delimitan su ventaja computacional.
Sin una recuperación relevante frente a L2 o una reducción temporal estable,
no se justifica anunciar una v1.5 sustancialmente superior ni agrandar el
manuscrito principal con toda esta campaña.

\section*{Reproducción}

Extraer la evidencia científica heredada de Online Resource 2 v1.4.2,
superponer sus escenarios prospectivos y luego este seguimiento. Desde
esa raíz ejecutar, sin Gurobi ni CPLEX:

\begin{verbatim}
python3 exploration/orf-v15-verified-20261019/check.py
python3 exploration/orf-v15-verified-20261019/check_selector.py
python3 exploration/orf-v15-verified-20261019/precision.py
python3 exploration/orf-v15-verified-20261019/analyze.py
python3 exploration/orf-v15-verified-20261019/check_statistics.py
python3 exploration/orf-v15-verified-20261019/projection_sensitivity.py
\end{verbatim}

Se retienen 6725 eventos de candidatos, sus vectores y todas las trazas.
VALIDATION.json, STATISTICS\_VALIDATION.json, SELECTOR\_VALIDATION.json y
PRECISION\_RESULTS.json documentan comprobaciones independientes y límites.
El DOI de la versión de envío permanece reservado para el final por
decisión del autor; este trabajo no publica ni envía un manuscrito.
\end{document}
